\documentclass[11pt, a4paper]{article}

\usepackage[T1]{fontenc}
\usepackage[utf8]{inputenc}
\usepackage{tgpagella}
\usepackage[spanish, es-noshorthands]{babel}
\usepackage{amsmath, amssymb, bm}
\usepackage{tabularx}
\usepackage[table, xcdraw]{xcolor}
\usepackage[most]{tcolorbox}
\usepackage[margin=2.5cm]{geometry}
\usepackage{ragged2e}
\usepackage{fancyhdr}

\pagestyle{fancy}
\fancyhf{}
\setlength{\headheight}{16pt}
\lhead{\textbf{UCB Sede Tarija} | Ing. Mecatrónica}
\chead{}
\rhead{IMT-121 Circuitos Electrónicos I | Checklist Lab A}
\lfoot{Prof: M.Sc. Ing. Bernardo Quiroga Turdera}
\rfoot{Página \thepage}
\renewcommand{\headrulewidth}{0.4pt}

\definecolor{ucbblue}{RGB}{0, 51, 102}

\begin{document}

\begin{tcolorbox}[colback=ucbblue!5!white, colframe=ucbblue, title=\textbf{Lista de Verificación: Cierre de Experiencia Integrada A}]
    \centering \textbf{Triage y Control de Trazabilidad: Teoría $\rightarrow$ SPICE $\rightarrow$ Banco}
\end{tcolorbox}

\textbf{Equipo:} \hrulefill \quad \textbf{Fecha:} \hrulefill

\vspace{0.3cm}
\noindent Marque con una "X" a medida que consolide la evidencia técnica. \textbf{Un montaje funcional no es una experiencia terminada}. Debe existir evidencia que conecte modelo, simulación y medición antes de desarmar el circuito.

\section*{1. Definición del Circuito y Modelo}
\begin{itemize}\setlength\itemsep{0.1em}
    \item[$\square$] Circuito base identificado correctamente en papel y simulador.
    \item[$\square$] Puerto $a-b$ definido claramente.
    \item[$\square$] Referencia de $V_a$ fijada respecto a una tierra común bien definida.
\end{itemize}

\section*{2. Superposición (Analítica, Simulada y Medida)}
\begin{itemize}\setlength\itemsep{0.1em}
    \item[$\square$] Cálculo teórico de contribuciones completado ($V_a^{(V1)} \approx 4.627\,\text{V}$ y $V_a^{(V2)} \approx 1.087\,\text{V}$).
    \item[$\square$] Simulación SPICE de los dos subcircuitos (fuentes individuales) completada.
    \item[$\square$] Medición experimental de $V_a^{(V1)}$ aislando físicamente $V_1$ (llevando $V_2 \rightarrow 0\,\text{V}$).
    \item[$\square$] Medición experimental de $V_a^{(V2)}$ aislando físicamente $V_2$ (llevando $V_1 \rightarrow 0\,\text{V}$).
    \item[$\square$] Comprobación de que la suma algebraica reproduce la respuesta total del circuito energizado.
\end{itemize}

\section*{3. Parámetros del Equivalente Thévenin / Norton}
\begin{itemize}\setlength\itemsep{0.1em}
    \item[$\square$] Medición de $V_{oc} = V_{th}$ realizada con el puerto $a-b$ \textbf{completamente abierto} (sin carga).
    \item[$\square$] Medición de $R_{th}$ realizada con la red \textbf{desenergizada} y medida mirando desde el puerto $a-b$.
    \item[$\square$] Valor analítico de la corriente de Norton reportado ($I_N = V_{th}/R_{th} \approx 2.628\,\text{mA}$).
    \item[$\square$] Dibujo esquemático explícito del equivalente de Thévenin y del equivalente de Norton.
\end{itemize}

\section*{4. Validación con Carga Externa}
\begin{itemize}\setlength\itemsep{0.1em}
    \item[$\square$] Respuesta medida ($V_L, I_L$) al conectar al menos una carga $R_L$ conocida al puerto $a-b$.
\end{itemize}

\section*{5. Interpretación y Discrepancias}
\begin{itemize}\setlength\itemsep{0.1em}
    \item[$\square$] Tabla comparativa (Teoría vs. SPICE vs. Medición) completamente llena.
    \item[$\square$] Cálculo de error relativo efectuado para las variables principales.
    \item[$\square$] Redacción de al menos \textbf{una hipótesis técnica} que justifique físicamente la discrepancia principal observada (ej. impedancia de instrumentos, tolerancia de componentes, falsos contactos, no usar tierra común).
\end{itemize}

\vspace{0.5cm}
\begin{center}
    \textbf{Firma del Instructor (Aprobación para Iniciar Exp. B):} \rule{6cm}{0.4pt}
\end{center}

\end{document}
